\subsection*{Texto C. Cuento: La silla}
\textit{Texto literario original de práctica.}

Cada tarde, Inés colocaba una silla frente a la puerta. No se sentaba: limpiaba el respaldo y regresaba al taller. Los vecinos suponían que esperaba a alguien. Cuando le preguntaban, respondía que el sol secaba mejor la madera a esa hora.

Una mañana llegó una carta sin remitente. Inés reconoció la letra de su hermano, ausente desde hacía doce años. La guardó bajo una caja de tornillos y siguió trabajando. Por la tarde llovió. Sacó la silla de todos modos y permaneció de pie a su lado, hasta que el agua le oscureció las mangas.

Al día siguiente reparó una de las patas, aunque no estaba rota. Después dejó la silla dentro y abrió la puerta. Un niño que pasaba quiso saber si ya no esperaba visitas. Inés sonrió: «Hoy pueden entrar». El niño miró el taller vacío, pero ella había vuelto a encender la lámpara.

\subsection*{Texto D. Ensayo literario: El borde de la página}
\textit{Texto literario original de práctica.}

Hay libros que recordamos por una frase y otros por el lugar donde los leímos. En mi ejemplar de viajes no conservo una sola fecha; sin embargo, una mancha de café me devuelve la mesa de una estación. La página no contiene aquel ruido, pero lo convoca.

Quien presta un libro suele pedir que se lo devuelvan intacto. Comprendo el cuidado: un objeto también merece respeto. Aun así, las marcas que deja una lectura no siempre son daños. Una anotación al margen puede ser la forma más breve de una conversación. El problema aparece cuando mi voz ocupa tanto espacio que ya no deja escuchar la del autor.

No propongo convertir cada volumen en un cuaderno de ejercicios. Prefiero pensar el margen como una orilla: permite acercarse al texto sin confundirlo con nosotros. Leer sería entonces aprender a habitar esa distancia. Hay ocasiones en que el mejor comentario es una línea; en otras, cerrar el lápiz y seguir leyendo.
